% EMNLP 2026 Main — ARR May 2026
% Generational Depth Contamination
\documentclass[11pt]{article}
\usepackage[review]{acl}
\usepackage{times}
\usepackage{latexsym}
\usepackage[T1]{fontenc}
\usepackage[utf8]{inputenc}
\usepackage{microtype}
\usepackage{graphicx}
\usepackage{amsmath}
\usepackage{amssymb}
\usepackage{booktabs}
\usepackage{multirow}
\usepackage{enumitem}
\usepackage{subcaption}
\usepackage[section]{placeins}

\newcommand{\Kstar}{{K^*}}
\newcommand{\Kempirical}{{K^*_{\text{empirical}}}}
\newcommand{\Ktheory}{{K^*_{\text{theory}}}}
\newcommand{\auc}{\text{AUC}}
\newcommand{\jsd}{\text{JSD}}

\title{Generational Depth Estimation of Recursive Synthetic Text:\\Measuring the Discrimination Boundary}

\author{Anonymous}

\begin{document}
\maketitle

% === Abstract ===
\begin{abstract}
Web-scale corpora are increasingly contaminated by recursive synthetic content, yet no method estimates a document's generational depth.
We measure a lower bound on the generational discrimination boundary, $\Kempirical$ (the maximum depth at which adjacent generations remain distinguishable), using 15 statistical features and an OBD-LightGBM classifier.
Across 10 model configurations (6 architecture families, 1B--8B), 3 source domains, and 4 decoding strategies, $\Kempirical \in \{1, 2, 3\}$ at $\theta{=}0.60$ (Cohen's $h{\approx}0.20$): general web text yields $\Kempirical{=}2$ in 6/10 configurations (2 of 6 families unanimously, plus a scale-dependent split in Pythia), while structured academic text (arXiv) yields $\Kempirical{=}3$, validated across two architectures.
A single generation pass suffices to render general web text indistinguishable from deeper synthetic content, a \emph{one-wash laundering} effect (first-generation text is detectable, but second-generation and beyond are not).
Model-free statistical features enable 80--100\% cross-model transfer with scorer-invariant feature rankings ($\rho > 0.90$); $\Kempirical$ agreement is scorer-independent for strong-signal configurations but scorer-sensitive at marginal boundaries (Qwen-continuation 1v2 $\auc = 0.611$, only 0.011 above $\theta$).
As a consistency check, domain-appropriate $\Kempirical$ correctly predicts that binary filtering outperforms graduated alternatives for web text when contamination causes detectable quality loss.
\end{abstract}

% === Section 1: Introduction ===
\section{Introduction}
\label{sec:introduction}

Web-scale text corpora are increasingly contaminated by synthetic content that has itself been generated from earlier synthetic content.
Estimates suggest that a substantial fraction of online text is machine-generated~\cite{sun2025aigt, spennemann2025quantifying}, and this share continues to rise.
When models train on the outputs of earlier models, the resulting multi-generational synthetic text degrades with each iteration: tail modes vanish, lexical diversity contracts, and surprisal distributions flatten~\cite{shumailov2024collapse, dohmatob2024model}.
This degradation arises in both \emph{retrain-chains} (iteratively training on model outputs) and \emph{prompt-chains} (recursively generating from previous outputs without retraining); we study the latter, which isolates generation-level information loss from training dynamics (Appendix~\ref{sec:extended_limitations}).
Yet the dominant paradigm treats all machine-generated text identically, discarding the generational depth dimension.

This gap cuts across three communities: binary detectors~\cite{basani2026diveye, zhang2024minkpp} produce no depth estimate (Section~\ref{sec:binary_confidence}); membership inference~\cite{zawalski2026codec} addresses a different question; the model collapse literature~\cite{shumailov2024collapse, alemohammad2024selfconsuming, schaeffer2025position} provides no per-text tool; and data selection methods~\cite{xie2023dsir, tirumala2023d4, drayson2025} ignore recursive depth.
No existing method estimates how many rounds of recursive generation a text has undergone.

The Data Processing Inequality guarantees that information degrades monotonically across generations, motivating a finite discrimination boundary $\Kstar$.
We empirically measure a lower bound $\Kempirical$ on the discrimination boundary (the maximum depth at which a classifier reliably distinguishes adjacent generations) using 15 statistical features and an Ordinal Binary Decomposition (OBD) LightGBM classifier.
We operationalize $\Kempirical$ as the deepest $k \geq 1$ where the pairwise $\auc_{k-1,k}$ exceeds $\theta = 0.60$ (Section~\ref{sec:formulation}).

Our contributions:
\begin{enumerate}[label=(\alph*),nosep]
  \item \textbf{First measurement of $\Kempirical$} across 10 model configurations (6 families, 1B--8B) and 4 decoding strategies at $\theta{=}0.60$ ($p \leq 0.002$; Section~\ref{sec:kstar_boundary}).
  6/10 configurations yield $\Kempirical{\geq}2$, spanning 2/6 families unanimously (Qwen, LLaMA) with a scale-dependent split in Pythia (1.4B: $\Kempirical{=}2$; 6.9B: $\Kempirical{=}1$); Qwen contributes 4/6 observations, with LLaMA-8B and Pythia-1.4B providing cross-family evidence.
  Both supervised OBD and zero-shot frontier LLMs fail comparably beyond $\Kempirical$ (McNemar $p > 0.05$; Section~\ref{sec:why_kstar}).
  \item \textbf{Domain dependence of $\Kempirical$}: general web text (C4, Wikipedia) yields $\Kempirical{=}2$, while structured academic text (arXiv) yields $\Kempirical{=}3$, consistent across model families (Section~\ref{sec:cross_domain}).
  \item \textbf{Model-free cross-model transfer}: 15 statistical features enable 80--100\% accuracy retention with scorer-invariant feature rankings ($\rho > 0.90$); $\Kempirical$ agreement holds across scorers for strong-signal configurations (Section~\ref{sec:transfer}).
\end{enumerate}

\noindent As a consistency check, $\Kempirical{=}2$ correctly predicts the strategy ordering for web-text filtering: binary curation minimizes degradation most effectively (MMLU $0.5941$ vs.\ base $0.6018$; Section~\ref{sec:filtering}).


% === Section 2: Related Work ===
\section{Related Work}
\label{sec:related_work}

\paragraph{Model collapse and recursive degradation.}
Iterative training on model outputs causes tail-distribution degradation~\cite{shumailov2024collapse, alemohammad2024selfconsuming, dohmatob2024model, dohmatob2025strong, fu2025theoretical}, diversity contraction~\cite{guo2024curious, padmakumar2024writing, briesch2023selfconsuming}, and predictable convergence~\cite{schaeffer2025position, kazdan2025collapse, gerstgrasser2024model, wang2025unveiling}.
Our prompt-chain operates in the replace regime~\cite{kazdan2025collapse}, making $\Kempirical = 2$ a conservative lower bound.
This literature characterizes macro dynamics but provides no per-text depth estimation, the gap we address.

\paragraph{Beyond binary detection.}
Binary detectors~\cite{mitchell2023detectgpt, bao2024fast, hans2024binoculars, basani2026diveye, zhang2024minkpp} achieve strong human-vs.-machine separation but produce no depth estimate.
Recent work extends beyond binary labels to collaboration degree~\cite{cheng2025beyond}, contamination hierarchies~\cite{mehta2025beyond}, and multigenerator benchmarks~\cite{wang2024semeval}; robustness concerns arise under paraphrasing~\cite{sadasivan2023detect, krishna2023paraphrasing} and cross-mode transfer~\cite{hao2025rewrite}; CoDeC~\cite{zawalski2026codec} addresses contamination via membership inference.
None estimates recursive generation depth.

\paragraph{Data selection and curation.}
DSIR~\cite{xie2023dsir}, D4~\cite{tirumala2023d4}, and DoReMi~\cite{xie2023doremi} are depth-agnostic; Drayson et al.~\cite{drayson2025} combined binary detection with importance resampling, a special case when $\Kstar = 1$.

\paragraph{Ordinal classification and probing.}
Our OBD classifier follows Frank \& Hall's~\cite{frank2001ordinal} binary decomposition; our neural probe extends probing~\cite{alain2017understanding, belinkov2022probing} to the discrimination boundary.


% === Section 3: Method ===
\section{Method}
\label{sec:method}

\begin{figure*}[t]
\centering
\includegraphics[width=\textwidth]{figures/pipeline.pdf}
\caption{Generational depth estimation pipeline. Seed documents are recursively passed through a generation operator $G_\delta$ to produce a depth chain. Statistical features are fed to the OBD-LightGBM classifier for depth prediction.}
\label{fig:pipeline}
\end{figure*}

\subsection{Problem Formulation}
\label{sec:formulation}

A generation operator $G_\delta$ (LLM + decoding strategy $\delta$) produces successive generations $X \to Y_1 \to \cdots \to Y_K$, where $Y_k = G_\delta(Y_{k-1})$ and $Y_0 \triangleq X$ (Figure~\ref{fig:pipeline}).
This sequence forms a Markov chain.
By the Data Processing Inequality, $I(X; Y_1) \geq I(X; Y_2) \geq \cdots \geq 0$: mutual information with the seed text degrades monotonically.
Together with Markov chain convergence, this provides heuristic motivation for expecting $\auc_{k-1,k}$ to decrease with $k$, though AUC and mutual information are not formally equivalent; the actual saturation depth must therefore be measured empirically.

We define the operational lower bound:
\begin{equation}
  \Kempirical = \max\!\bigl\{k \!\geq\! 1 : \auc_{j-1,j} > \theta,\;\forall\, j \!\leq\! k\bigr\},
  \label{eq:kstar}
\end{equation}
which requires all pairs up to depth $k$ to exceed the threshold (the $\forall j$ constraint is naturally satisfied because $\auc_{k-1,k}$ decreases monotonically above $\theta$ across all tested configurations; Table~\ref{tab:kstar_ci}).
The theoretical upper bound $\Ktheory$ (the true maximum discriminable depth given optimal features) is not directly estimated in this work; $\Kempirical$ provides a feature-set-specific lower bound on $\Ktheory$.
Here $\auc_{k-1,k}$ is the point-estimate pairwise AUC and $\theta = 0.60$ is a small-effect-size threshold (Cohen's $h \approx 0.20$\footnote{$h = 2\arcsin\!\sqrt{0.60} - 2\arcsin\!\sqrt{0.50} \approx 0.2014$, a ``small'' effect size by Cohen's (1988) convention.}; justified below).
Bootstrap 95\% CIs (1{,}000 resamples) provide uncertainty quantification: we verify that boundary-pair CI lower bounds exceed $\theta$ and non-boundary-pair CI upper bounds fall below $\theta$ (Section~\ref{sec:kstar_boundary}).

\paragraph{Threshold justification.}
$\theta{=}0.60$ is supported by three independent lines of evidence.
First, it corresponds to Cohen's $h \approx 0.20$, a conventional small-effect-size boundary.
Second, a 10{,}000-iteration permutation test yields a 95th-percentile null of ${<}\,0.55$ and $p < 0.002$ for all boundary comparisons, placing $\theta{=}0.60$ well above the empirical null.
Third, as post-hoc confirmatory evidence, the conclusion $\Kempirical{=}2$ is insensitive to the exact threshold: 1v2 AUC values cluster at $0.61$--$0.70$ while 2v3 values cluster at $0.55$--$0.59$, creating a ${\sim}$5--11\,pp discontinuity that any $\theta \in [0.55, 0.65]$ would separate (Figure~\ref{fig:kstar_theta}; Appendix~\ref{sec:extended_limitations}).
Only at extremes does $\Kempirical$ shift: $\theta{=}0.55$ yields $\Kempirical{=}3$ for two Qwen configurations; $\theta{=}0.70$ yields $\Kempirical{=}1$ across all configurations.
The downstream filtering evaluation (Section~\ref{sec:filtering}) uses the OBD classifier without re-optimizing $\theta$, avoiding circularity.

\subsection{Features and Classifier}
\label{sec:features}

Each document is characterized by a 15-dimensional feature vector spanning four groups: \emph{surprisal moments} (mean, std, skewness, kurtosis; 4~dims), \emph{surprisal derivatives} (mean and std of first- and second-order differences, plus first-order skewness; 5~dims), \emph{vocabulary diversity} (type-token ratio, hapax legomena ratio, self-BLEU; 3~dims), and \emph{frequency tail statistics} (frequency kurtosis, entropy, rare token fraction; 3~dims).
All features are computed from token-level surprisal values via an external scorer model (full definitions in Table~\ref{tab:s5_features}, Appendix).
Self-BLEU is computed within each CV fold using only same-fold, same-depth reference samples to prevent leakage; it ranks \#1 in 1/8 conditions (greedy decoding) but is absent from the top-3 in all 7 others.
Removing self-BLEU reduces mean $\auc_{1v2}$ by $< 1$\,pp across all models (14D $\Kempirical = 2$ for all 4 configurations); in deployment where depth labels are unavailable, self-BLEU can be replaced with corpus-level self-similarity or omitted with minimal performance loss.
Hand-crafted features enable cross-generator transfer without retraining and permit attribution of discriminative power to specific text properties.

We use \emph{Ordinal Binary Decomposition} (OBD), decomposing the $(K{+}1)$-class problem into $K$ cumulative binary sub-classifiers, each estimating $P(\text{depth} > k \mid \mathbf{x})$ via LightGBM (hyperparameters in Table~\ref{tab:s6_lgbm_params}).
OBD achieves lower MAE than multiclass LightGBM (e.g., 0.984 vs.\ 1.016 on Qwen-rewrite) while maintaining superior pairwise $\auc$ (mean 0.6095 vs.\ 0.4883 for ordinal logistic regression).

\subsection{Experimental Setup}
\label{sec:experimental_setup}

We evaluate six configurations spanning three model families (Table~\ref{tab:setup}): Qwen2.5-1.5B-Instruct (rewrite and continuation modes), Qwen2.5-1.5B (base, continuation), Qwen2.5-7B (continuation), Pythia-1.4B and OLMo-1B (continuation).
Each uses 5{,}000 C4 seeds truncated to 256 tokens, recursively generated for $K{=}5$ iterations (30{,}000 documents per configuration).
Default decoding: nucleus $p{=}0.95$, $T{=}1.0$; decoding ablation covers greedy, $T{=}0.7$, and $T{=}1.2$.
Continuation mode truncates seeds to 128 tokens and extends to 256; rewrite mode paraphrases the full text.
Evaluation: 5-fold GroupKFold cross-validation grouped by seed document; pairwise $\auc$ as primary metric.

\begin{table}[t]
\centering
\small
\caption{Experimental configurations for depth-chain generation. \textbf{Config}: shorthand identifier; \textbf{Model}: architecture and parameter count; \textbf{Mode}: generation strategy (continuation or rewrite); $N$: number of seed documents per depth level.}
\label{tab:setup}
\begin{tabular}{llcc}
\toprule
\textbf{Config} & \textbf{Model} & \textbf{Mode} & $N$ \\
\midrule
Qwen-rw   & Qwen2.5-1.5B-Instruct & Rewrite & 5{,}000 \\
Qwen-cont & Qwen2.5-1.5B-Instruct & Continuation & 5{,}000 \\
Qwen-base & Qwen2.5-1.5B          & Continuation & 5{,}000 \\
Pythia    & Pythia-1.4B            & Continuation & 5{,}000 \\
OLMo      & OLMo-1B               & Continuation & 5{,}000 \\
Qwen-7B   & Qwen2.5-7B            & Continuation & 5{,}000 \\
\bottomrule
\end{tabular}
\end{table}


% === Section 4: Experiments ===
\section{Experiments}
\label{sec:experiments}

We evaluate generational depth estimation on ten model configurations spanning six architecture families:
Pythia-1.4B, Pythia-6.9B, OLMo-1B, Qwen2.5-1.5B-Instruct (rewrite and continuation), Qwen2.5-1.5B (base), Qwen2.5-7B, LLaMA-3.1-8B, Mistral-7B, and Gemma-2-2B, all in continuation mode except where noted.
All experiments use the OBD-LightGBM classifier with the 15-dimensional feature set (Section~\ref{sec:method}), with 5{,}000 texts per depth level (generations 0--5).

%% ====================================================================
\subsection{$\Kempirical \leq 2$ Across Model Configurations (C4)}
\label{sec:kstar_boundary}

\begin{table*}[t]
\centering
\small
\caption{Pairwise $\auc$ for adjacent depth pairs with bootstrap 95\% CIs and permutation $p$-values for boundary rows.
$\Kempirical = 2$ requires 1v2 $\auc > \theta{=}0.60$ and 2v3 $\auc < \theta$; CIs quantify boundary uncertainty.
Extended results (all models) in Table~\ref{tab:kstar_ci_extended}.}
\label{tab:kstar_ci}
\begin{tabular}{lcccccccccc}
\toprule
 & \multicolumn{2}{c}{Qwen-rw} & \multicolumn{2}{c}{Pythia} & \multicolumn{2}{c}{OLMo} & \multicolumn{2}{c}{Qwen-cont} & \multicolumn{2}{c}{Qwen-7B\rlap{$^\ddagger$}} \\
\cmidrule(lr){2-3}\cmidrule(lr){4-5}\cmidrule(lr){6-7}\cmidrule(lr){8-9}\cmidrule(lr){10-11}
 & $\auc$ & 95\% CI & $\auc$ & 95\% CI & $\auc$ & 95\% CI & $\auc$ & 95\% CI & $\auc$ & 95\% CI \\
\midrule
0v1 & .9956 & & .9932 & & .9860 & & .9971 & & .9638 & [.961,\,.967] \\
1v2 & .6955 & [.686,\,.704]\rlap{$^*$} & .6539 & [.645,\,.663]\rlap{$^*$} & .5080 & [.496,\,.519] & .6112 & [.602,\,.621]\rlap{$^*$} & .6168 & [.608,\,.626]\rlap{$^*$} \\
2v3 & .5893 & [.580,\,.599] & .5514 & [.542,\,.563] & .4973 & [.486,\,.509] & .5609 & [.551,\,.571] & .5492 & [.539,\,.559] \\
3v4 & .5435 & & .5152 & & .5024 & & .5409 & & .5276 & \\
4v5 & .5340 & & .5136 & & .4942 & & .5141 & & .5085 & \\
\bottomrule
\multicolumn{11}{l}{\footnotesize $^*$Permutation $p < 0.001$ for all 1v2 comparisons shown.} \\
\multicolumn{11}{l}{\footnotesize $^\ddagger$Generator: Qwen2.5-7B (fp16); scorer: Qwen2.5-1.5B.} \\
\multicolumn{11}{l}{\footnotesize Full results for all 10 configurations (including Qwen-7B, Qwen-base, Gemma-2-2B, and Pythia-6.9B) are provided in Table~\ref{tab:kstar_ci_extended}.}
\end{tabular}
\end{table*}

\begin{figure}[t]
\centering
\includegraphics[width=\columnwidth]{figures/fig2_auc_decay.pdf}
\caption{Pairwise $\auc$ decay across depth pairs for model-mode configurations (solid) and decoding strategies (dashed); shaded bands show bootstrap 95\% confidence intervals. The horizontal dashed line marks $\theta = 0.60$.}
\label{fig:auc_decay}
\end{figure}

Table~\ref{tab:kstar_ci} and Figure~\ref{fig:auc_decay} show $\Kempirical \leq 2$ across all 10 model configurations: 6/10 configurations yield $\Kempirical{=}2$ (2/6 families unanimously: Qwen with 4 configs, LLaMA-8B; Pythia-1.4B provides additional cross-family evidence) (1v2 $\auc$ range: 0.6112--0.6955; all CI lower bounds $> \theta$; permutation $p < 0.001$), while OLMo-1B, Mistral-7B, Pythia-6.9B, and Gemma-2-2B yield $\Kempirical{=}1$ (1v2 $\auc$: 0.508, 0.594, 0.596, and 0.598, all below $\theta$; Section~\ref{sec:boundary_cases}).
LLaMA-3.1-8B confirms $\Kempirical{=}2$ (0v1 $\auc = 0.9225$, 1v2 $\auc = 0.6123\;[0.603,\,0.622]$; CI lower bound above $\theta$; Table~\ref{tab:kstar_ci_extended}).
All four decoding strategies likewise yield $\Kempirical{=}2$ (Table~\ref{tab:s2_decoding}).
The 1v2 spread due to generation mode ($\Delta = 8.4$\,pp within Qwen) exceeds the spread across $\Kempirical{=}2$ continuation models ($\Delta = 4.3$\,pp).

\paragraph{Scale and instruction-tuning invariance.}
$\Kempirical{=}2$ persists from 1.5B through 8B: Qwen-7B (1v2 $\auc = 0.6168$ $[0.608, 0.626]$) and LLaMA-3.1-8B (1v2 $\auc = 0.6123$) both exceed $\theta$, despite progressively reduced 0v1 signal (0.9638 and 0.9225 vs.\ $> 0.99$ at $\leq$1.5B).
A base-vs.-instruct ablation at 1.5B likewise yields $\Kempirical = 2$ (1v2 $\auc = 0.626$), confirming the boundary is not an instruction-tuning artifact.

\begin{figure}[t]
\centering
\includegraphics[width=\columnwidth]{figures/fig_kstar_theta.pdf}
\caption{$\Kempirical(\theta)$ sensitivity curve for six configurations. $\Kempirical$: maximum discriminable depth; $\theta$: AUC decision threshold. Five of six configurations yield $\Kempirical{=}2$ at $\theta{=}0.60$ (dashed line); OLMo yields $\Kempirical{=}1$. At $\theta{=}0.70$, all drop to $\Kempirical{=}1$.}
\label{fig:kstar_theta}
\end{figure}

Figure~\ref{fig:kstar_theta} shows $\Kempirical$ as a function of $\theta$: five of six configurations yield $\Kempirical = 2$ within $\theta \in [0.51, 0.61]$ (OLMo: $\Kempirical = 1$ throughout); at $\theta{=}0.55$, two Qwen configurations reach $\Kempirical = 3$; at $\theta{=}0.70$, all drop to $\Kempirical = 1$.
All four decoding strategies likewise yield $\Kempirical = 2$ at $\theta{=}0.60$ (1v2 $\auc \in [0.687, 0.704]$; Table~\ref{tab:s2_decoding}).

%% ====================================================================
\subsection{Cross-Model Transfer}
\label{sec:transfer}

\begin{figure}[t]
\centering
\includegraphics[width=0.85\columnwidth]{figures/fig3_transfer_heatmap.pdf}
\caption{Cross-model transfer heatmap. Each cell shows 6-class ordinal accuracy when trained on the row model and tested on the column model, divided by in-domain accuracy (retention rate). Pythia$\leftrightarrow$OLMo: 90--92\%; Qwen-rewrite$\to$continuation: 61--63\%.}
\label{fig:transfer_heatmap}
\end{figure}

Continuation-mode classifiers transfer well across generators (Figure~\ref{fig:transfer_heatmap}; full matrix in Table~\ref{tab:s1_transfer}): Pythia$\leftrightarrow$OLMo retains 90--92\%, Qwen-cont$\leftrightarrow$Pythia 84--85\%, Qwen-cont$\leftrightarrow$OLMo 80--85\% (0v1 cross-model $\auc > 0.96$ for all pairs).
Qwen-rewrite classifiers transfer poorly (61--63\%), confirming reliance on mode-specific artifacts rather than generalizable depth features.

\paragraph{Cross-scorer validation.}
Replacing the self-scorer with Pythia-1.4B or OLMo-1B preserves feature rankings ($\rho > 0.90$) but attenuates absolute discriminability at the boundary (Table~\ref{tab:crossscorer}).
Feature ranking independence is a strong result: the same features dominate under all scorers.
$\Kempirical$ agreement, however, is signal-strength dependent: Qwen-rewrite (strong signal, 1v2 $\auc = 0.696$) retains $\Kempirical{=}2$ under all scorers, but Qwen-continuation (marginal, 1v2 $\auc = 0.611$, only 0.011 above $\theta$) drops to $\Kempirical{=}1$ (Pythia: $0.5923$; OLMo: $0.5858$).
A three-way ANOVA (Generator $\times$ Scorer $\times$ Depth-Pair; Table~\ref{tab:anova_crossscorer}) confirms a statistically significant but modest Scorer main effect ($\eta^2_p{=}.166$, $p < .001$) and a significant Generator$\times$Scorer interaction ($\eta^2_p{=}.146$), indicating that absolute AUC values are scorer-dependent even when $\Kempirical$ classifications largely agree.

%% ====================================================================
\subsection{Feature Ordering and Generation Mode}
\label{sec:feature_degradation}
\label{sec:generation_mode}

Among continuation-mode models, single-feature $\auc$ rankings are highly correlated ($\rho \geq 0.85$; Table~\ref{tab:top_features}).
The top-3 features (\texttt{surp\_d1\_std}, \texttt{surp\_std}, \texttt{surp\_d2\_std}) are shared across all three families, revealing that autoregressive generation progressively damps surprisal variability.
Qwen-rewrite diverges (\texttt{freq\_entropy} dominates), consistent with its poor cross-model transfer (61--63\%).
Within Qwen, rewrite yields 1v2 $\auc = 0.6955$ vs.\ continuation's $0.6112$ ($\Delta = +8.4$\,pp); continuation mode is therefore recommended as the canonical condition.

%% ====================================================================
\subsection{Neural Probe}
\label{sec:neural_probe}

Replacing the 15D features with frozen encoder hidden states (MLP probe; Table~\ref{tab:neural_probe_crossmodel}), Qwen achieves $\Kempirical^{\text{neural}} = 3$ (2v3 $\auc = 0.6159$), while Pythia remains at $\Kempirical^{\text{neural}} = 2$ ($\auc_{2v3} = 0.5963$); OLMo yields $\Kempirical^{\text{neural}} = 1$ (1v2 $\auc = 0.5054$, consistent with its 15D boundary).
The Qwen-only elevation reflects self-familiarity artifacts that transfer poorly (Section~\ref{sec:transfer}); the 15D features yield $\Kempirical = 2$ for most families (OLMo: $\Kempirical = 1$).

\begin{minipage}{\columnwidth}
\centering
\footnotesize
\captionof{table}{Neural probe pairwise AUC using frozen-encoder MLP. Each cell reports AUC for classifying depth $d$ vs.\ $d{+}1$ text. $\Kempirical^{\text{neural}}$: maximum depth where AUC exceeds $\theta = 0.60$.}
\label{tab:neural_probe_crossmodel}
\begin{tabular}{lccc}
\toprule
\textbf{Model} & \textbf{2v3 $\auc$} & \textbf{95\% CI} & $\Kempirical^{\text{neural}}$ \\
\midrule
Qwen-1.5B  & 0.6159 & $[0.610,\,0.622]$ & 3 \\
Pythia-1.4B & 0.5963 & $[0.589,\,0.603]$ & 2 \\
OLMo-1B    & 0.5056 & $[0.497,\,0.516]$ & 1\rlap{$^\dagger$} \\
\bottomrule
\multicolumn{4}{l}{\footnotesize $^\dagger$OLMo 1v2 $\auc = 0.5054$ (below $\theta$); $\Kempirical^{\text{neural}} = 1$.}
\end{tabular}
\end{minipage}

%% ====================================================================
\subsection{Cross-Domain Variation}
\label{sec:cross_domain}

The preceding experiments use C4 (general web text).
To test domain dependence, we evaluate on Wikipedia ($n{=}5{,}000$) and arXiv abstracts ($n{=}2{,}200$) using two model families.
Wikipedia yields $\Kempirical{=}2$ (1v2 $\auc = 0.623$, 2v3 $\auc = 0.571$), matching the C4 baseline.
Both Qwen-1.5B and Pythia-1.4B yield $\Kempirical{=}3$ for arXiv (Table~\ref{tab:domain_kstar}), demonstrating that structured academic text permits discrimination to greater depth independently of architecture.
The elevated signal reflects higher per-generation distributional shift (1v2 JSD $= 0.077$ vs.\ ${\sim}0.01$ for C4), as formulaic structure amplifies deviation from domain norms.

\begin{minipage}{\columnwidth}
\centering
\footnotesize
\captionof{table}{Cross-domain $\Kempirical$ comparison. Rows: generator model and text domain (C4, Wikipedia, arXiv). $\Kempirical{=}3$ for arXiv indicates structured text preserves discriminability to greater depth.}
\label{tab:domain_kstar}
\begin{tabular}{llccc}
\toprule
\textbf{Domain} & \textbf{Model} & \textbf{1v2 $\auc$} & \textbf{2v3 $\auc$} & $\Kempirical$ \\
\midrule
arXiv & Qwen-1.5B  & 0.858 & 0.6436 & 3 \\
arXiv & Pythia-1.4B & 0.805 & 0.6545 & 3 \\
C4    & Qwen-1.5B  & 0.63  & 0.56   & 2 \\
Wikipedia & Qwen-1.5B & 0.62 & 0.55 & 2 \\
\bottomrule
\end{tabular}
\end{minipage}

%% ====================================================================
\subsection{Binary Detector Confidence}
\label{sec:binary_confidence}

Binary detectors (DivEye~\cite{basani2026diveye}, Binoculars~\cite{hans2024binoculars}) show moderate rank correlations with depth ($|\rho| = 0.06$--$0.46$; Table~\ref{tab:s9_binary_spearman}) but fail as depth proxies: 6-class accuracy is at or below chance ($\leq 0.199$ vs.\ 0.167 baseline), pairwise $\auc$ yields $\Kempirical^{\text{binary}} \leq 1$, and multiple conditions exhibit non-monotonic AUC at deeper pairs.
The observed correlation is an artifact of 0v1 separation rather than fine-grained depth signal.


% === Section 5: K*-Aware Data Filtering ===
\section{Consistency Check: $\Kstar$-Guided Filtering}
\label{sec:filtering}

$\Kempirical = 2$ makes a falsifiable prediction: binary curation should suffice, and graduated performance should converge toward binary as the decay steepness increases.
We test this as a consistency check on the boundary measurement, not as an independent contribution.
We compare six strategies on 30{,}000 documents (depths 0--5): (i)~\textbf{no-filter}, (ii)~\textbf{binary-ours} ($w = 1 - P(\text{synthetic})$ from 15D LightGBM), (iii)~\textbf{binary-Drayson}~\citep{drayson2025}, (iv)~\textbf{graduated} ($w = \exp(-\alpha \cdot \hat{d})$, $\alpha \in \{0.1,\ldots,0.9\}$), (v)~\textbf{surprisal} (retain above median), and (vi)~\textbf{DSIR}~\citep{xie2023dsir}.
We fine-tune Qwen2.5-1.5B-Instruct via LoRA (rank~8, $\alpha_{\text{LoRA}}{=}16$) for 3 epochs on MMLU and HellaSwag across 5 seeds (51 total runs; this 100\% synthetic setting is an upper bound~\cite{shumailov2024collapse, dohmatob2024model}).

\begin{table}[t]
\centering
\small
\caption{Downstream task performance (MMLU, HellaSwag) under six data-filtering strategies. Binary-ours uses the 15D LightGBM depth classifier; Graduated applies exponentially decaying weights $w = \exp(-\alpha \hat{d})$. Mean $\pm$ std over 5 seeds. Full $\alpha$ sweep in Appendix~\ref{sec:extended_limitations}.}
\label{tab:tab_filtering_main}
\begin{tabular}{lcc}
\toprule
\textbf{Strategy} & \textbf{MMLU} & \textbf{HellaSwag} \\
\midrule
Base model (no training) & 0.6018 & 0.6821 \\
\midrule
Binary-ours ($\Kstar{=}2$) & \textbf{0.5941}$\pm$0.0008 & \textbf{0.6701}$\pm$0.0006 \\
Binary-Drayson & 0.5913$\pm$0.0009 & 0.6691$\pm$0.0004 \\
Surprisal & 0.5894$\pm$0.0009 & 0.6641$\pm$0.0006 \\
Graduated ($\alpha{=}0.9$) & 0.5892$\pm$0.0008 & 0.6701$\pm$0.0007 \\
DSIR & 0.5869$\pm$0.0003 & 0.6631$\pm$0.0007 \\
No filter & 0.5865$\pm$0.0010 & 0.6638$\pm$0.0008 \\
\bottomrule
\end{tabular}

\end{table}

Binary-ours minimizes MMLU degradation most effectively (Table~\ref{tab:tab_filtering_main}), reducing the gap from $-$1.53\,pp (no-filter) to $-$0.77\,pp (${\sim}$50\% reduction; $p_{\text{corr}} < 0.05$ vs.\ all alternatives).
These effects are measured in a 100\% synthetic upper-bound setting; real-world contamination ratios are typically lower, and effect sizes would scale accordingly.
The strategy ordering is interpretable through $\Kempirical$: binary filtering exploits the one reliable boundary (0v1), while graduated strategies assign weights to depths 2--5 based on near-random classifier outputs, injecting noise.
Graduated filtering converges monotonically toward binary: at $\alpha{=}0.9$, it matches binary-ours on HellaSwag exactly ($\Delta = 0$, $p_{\text{corr}} = 1.00$) but remains significantly worse on MMLU ($p_{\text{corr}} = 0.011$).
DSIR is indistinguishable from no-filter, confirming that depth-aware features are necessary for curation in recursively contaminated corpora (Figure~\ref{fig:filtering_strategies} in Appendix).
At realistic contamination ratios ($\leq$50\%), all filtering strategies converge ($p > 0.12$), consistent with the one-wash laundering finding that depth-aware curation offers marginal benefit when $\Kempirical{=}2$ (Appendix~\ref{sec:dose_response}).


% === Section 6: Discussion ===
\section{Discussion}
\label{sec:discussion}
\label{sec:why_kstar}
\label{sec:one_wash}
\label{sec:negative_graduated}
\label{sec:repr_dependence}
\label{sec:boundary_cases}
\label{sec:length_confound}

\paragraph{One-wash laundering (web text).}
For general web text, $\Kempirical{=}2$ reflects concentrated information loss: 0v1 is near-trivially separable ($\auc \geq 0.92$), 1v2 drops by ${\sim}$30\,pp, and 2v3 falls below $\theta$.
A single generation pass renders text indistinguishable from multiply recycled content, limiting depth-aware curation to a binary boundary.
Converging evidence supports a data-intrinsic interpretation: OBD achieves 96\% depth-0 recall but only 14--22\% for depths 2--4; frontier LLMs show 8--38\% (McNemar $p > 0.05$ vs.\ OBD; Table~\ref{tab:s3_llm_zeroshot}).

\paragraph{Graduated filtering failure.}
For web text, 2v3 $\auc < \theta$ across all model configurations, so depth assignments for $d \geq 2$ approximate random labels; graduated weighting injects noise, explaining convergence toward binary at high $\alpha$ (Appendix~\ref{sec:extended_limitations}).
DSIR fails differently: bigram hashing captures domain shift but not depth-correlated signatures.

\paragraph{Representation dependence.}
$\Kempirical$ is representation-dependent: frozen Qwen embeddings yield $\Kempirical^{\text{neural}} = 3$, while Pythia yields $\Kempirical^{\text{neural}} = 2$ and OLMo yields $\Kempirical^{\text{neural}} = 1$, consistent with its 15D boundary (Table~\ref{tab:neural_probe_crossmodel}).
The 15D features trade one depth level of sensitivity for cross-model consistency (6/10 configurations $\Kempirical{=}2$), making them the preferred operational choice.
Feature rankings are scorer-independent ($\rho > 0.90$); $\Kempirical$ agreement is scorer-independent for strong-signal configurations (e.g., Qwen-rewrite retains $\Kempirical{=}2$ under all scorers) but scorer-sensitive at marginal boundaries (Qwen-continuation drops to $\Kempirical{=}1$; length confound ruled out in Appendix~\ref{sec:extended_limitations}).
ANOVA confirms this layered pattern: Scorer explains modest variance ($\eta^2_p{=}.166$) with a significant Generator$\times$Scorer interaction ($\eta^2_p{=}.146$; Table~\ref{tab:anova_crossscorer}).

\paragraph{Domain dependence of $\Kempirical$.}
Web text (C4, Wikipedia) consistently yields $\Kempirical{=}2$, while structured academic text (arXiv) yields $\Kempirical{=}3$ across two independent model families (Qwen-1.5B, Pythia-1.4B; Section~\ref{sec:cross_domain}).
Domain-specific stylistic constraints amplify per-generation distributional shift, making contamination in structured corpora trackable to greater depth.

\paragraph{Prompt-chain vs.\ retrain-chain $\Kempirical$.}
Retrain-chain experiments (Pythia-1.4B LoRA) yield $\Kempirical{\geq}4$ (1v2 $\auc = 0.9484$, 2v3 $= 0.8888$, 3v4 $= 0.7982$), far exceeding the prompt-chain boundary of $\Kempirical{=}2$.
The elevated retrain-chain signal arises from cumulative model-adaptation artifacts (vocabulary narrowing, repetition bias) rather than text-intrinsic degradation.
Prompt-chain $\Kempirical{=}2$ therefore captures the text-level discrimination boundary relevant to data curation; the separation between the two protocols validates our experimental design choice.

\paragraph{Boundary cases and scale effect.}
Of 10 configurations on C4, 6 yield $\Kempirical{=}2$ (2/6 families unanimously: Qwen contributes 4 observations, LLaMA-8B one; Pythia-1.4B adds cross-family evidence but Pythia-6.9B yields $\Kempirical{=}1$); OLMo-1B ($\Kempirical{=}1$; 1v2 $\auc = 0.508$, CI $[0.496, 0.519]$), Mistral-7B ($\Kempirical{=}1$; 1v2 $\auc = 0.594$, CI $[0.584, 0.603]$), Pythia-6.9B ($\Kempirical{=}1$; 1v2 $\auc = 0.596$), and Gemma-2-2B ($\Kempirical{=}1$; 1v2 $\auc = 0.598$, CI $[0.588, 0.608]$) fall below $\theta{=}0.60$, validating $\theta$ as an effective discriminator (gap: $0.508$--$0.594$ vs.\ $0.611$--$0.696$).
Larger models produce more human-like text, reducing 0v1 JSD: 0.15--0.27 at $\leq$1.5B, 0.077 at 8B (LLaMA-3.1-8B).
This trend suggests that models at 13B+ may increasingly yield $\Kempirical{=}1$; this speculative prediction warrants future verification.

\paragraph{Decoding strategy as confound.}
Pythia-6.9B under greedy decoding yields $\Kempirical{=}2$ (1v2 $\auc = 0.854$), but switching to nucleus sampling ($p{=}0.95$) reduces the boundary to $\Kempirical{=}1$ (1v2 $\auc = 0.596$), providing direct evidence that greedy repetition artifacts inflate depth separability.
The borderline 1v2 score (2 of 5 folds exceed $\theta{=}0.60$) illustrates that the $\Kempirical$ boundary is continuous rather than a binary cliff.
Pythia exhibits scale-dependent behavior: Pythia-1.4B yields $\Kempirical{=}2$ (1v2 $\auc = 0.654$) while Pythia-6.9B yields $\Kempirical{=}1$ (1v2 $\auc = 0.596$), confirming that increased model scale does not monotonically raise $\Kempirical$.


% === Section 7: Conclusion ===
\section{Conclusion}
\label{sec:conclusion}

We measure a lower bound on the generational depth discrimination boundary $\Kempirical$ for recursive synthetic text.
$\Kempirical \in \{1, 2, 3\}$ depends on source domain and model characteristics: general web text yields $\Kempirical{=}2$ in 6/10 configurations (2/6 families unanimously, with a scale-dependent split in Pythia; 1B--8B; $p \leq 0.002$), with $\Kempirical{=}3$ for structured academic text.
This one-wash laundering effect means that a single generation pass renders web text indistinguishable from deeper recursive content, limiting depth-aware curation to a binary boundary.
For practitioners, binary filtering (depth 0 vs.\ $\geq$1) captures the only reliable discrimination signal; graduated weighting beyond $\Kempirical$ adds noise.
$\Kempirical$ is feature-set dependent (15D provides a conservative lower bound; neural probes extend it by one level for some architectures) and established for 1B--8B models; scaling beyond 8B, additional structured domains, and retrain-chain $\Kstar$ remain open.


% === Limitations (ACL required, unnumbered) ===
\clearpage
\section*{Limitations}

Our findings are bounded by:
(1)~\emph{Scope}: model scale $\leq$8B (6 architecture families), English text (C4, Wikipedia, arXiv), four decoding strategies; arXiv $\Kempirical{=}3$ is confirmed with two model families (Qwen-1.5B, Pythia-1.4B) but with reduced sample size ($n{=}2{,}200$ per model).
(2)~\emph{Protocol}: prompt-chain rather than retrain-chain. Retrain-chain validation (Pythia-1.4B LoRA, Appendix~\ref{sec:extended_limitations}) yields $\Kempirical{\geq}4$ due to cumulative model-adaptation artifacts; prompt-chain isolates text-intrinsic degradation and provides the operationally relevant boundary for data curation.
(3)~\emph{Threshold}: $\Kempirical$ ranges from 2--3 at $\theta{=}0.55$ to 1 at $\theta{=}0.70$; 6/10 configurations yield $\Kempirical{=}2$ at $\theta{=}0.60$ (OLMo-1B, Mistral-7B, Gemma-2-2B, Pythia-6.9B: $\Kempirical{=}1$).
(4)~\emph{Marginal boundary}: Qwen-continuation's 1v2 CI lower bound exceeds $\theta$ by only 0.002; cross-scorer validation drops this to $\Kempirical{=}1$ (Section~\ref{sec:transfer}).
(5)~\emph{Feature set}: the 15D features yield a feature-set-specific lower bound ($\Ktheory$ may be strictly higher); neural probe $\Kempirical$ is model-dependent ($3$ for Qwen, $2$ for Pythia, $1$ for OLMo).
(6)~\emph{Segment length}: $\Kempirical{=}2$ confirmed for 128 and 256 tokens (CI lower $> \theta$); 512 tokens yields point estimate $\Kempirical{=}2$ but CI lower marginally below $\theta$ ($0.5979$), indicating the boundary weakens at longer generation lengths.
(7)~\emph{Curation applicability}: the practical benefit of $\Kstar$-aware curation depends on contamination severity; when synthetic data proportion is low enough to leave overall training quality intact, filtering strategies converge (Appendix~\ref{sec:dose_response}).
Our filtering validation (Section~\ref{sec:filtering}) was conducted exclusively on Qwen-1.5B continuation-mode text; extending to $\Kempirical{=}1$ generators remains future work.
Seed robustness testing (Appendix~\ref{sec:seed_robustness}) shows $\Kempirical$ is stable under seed variation for strong-signal configurations, but marginal boundaries are seed-sensitive (absolute AUC shifts up to ${\pm}9$\,pp).
Additional limitations (cross-model scoring sensitivity, feature ranking variability at the 1v2 boundary, structural composition confound analysis) appear in Appendix~\ref{sec:extended_limitations}.

% === Ethics Statement (ACL required, unnumbered) ===
\section*{Ethics Statement}

Generational depth estimation is intended for data curation and corpus auditing, positive-sum applications that improve training data quality.
Our experiments use the publicly available C4 corpus and do not involve personally identifiable information.
A potential dual-use concern is that depth estimation techniques could inform adversarial strategies to evade detection (e.g., by calibrating rewriting to mimic earlier-generation statistics); we note that this risk is shared by all detection research and that the binary detection boundary ($\auc \geq 0.96$ at 0v1) remains robust.

\bibliography{references}

\appendix
\section{Supplementary Figures}

\begin{figure}[t]
\centering
\includegraphics[width=\columnwidth]{figures/fig_filtering_strategies.pdf}
\caption{Downstream performance of all filtering strategies (5 seeds, $\pm$1 std). Binary-ours (green) minimizes MMLU degradation; graduated $\alpha$ sweep shows monotonic convergence toward binary. Dashed line: base model (no fine-tuning).}
\label{fig:filtering_strategies}
\end{figure}

%% [Merged into Figure 2] \begin{figure}[t]
%% [Merged into Figure 2]   \centering
%% [Merged into Figure 2]   \includegraphics[width=\columnwidth]{figures/fig_s1_auc_decay_ci.pdf}
%% [Merged into Figure 2]   \caption{Pairwise $\auc$ decay curves for five model-mode configurations with bootstrap 95\% confidence intervals. Error bars computed from 1{,}000 bootstrap resamples. The shaded region marks the $\Kempirical=2$ boundary where $\auc$ drops below $\theta=0.60$. This figure complements Figure~\ref{fig:auc_decay} in the main text by showing statistical uncertainty for each measurement point.}
%% [Merged into Figure 2]   \label{fig:s1_auc_decay_ci}
%% [Merged into Figure 2] \end{figure}

\section{Supplementary Tables}

%%% Feature ranking table (moved from main text)
\begin{table}[t]
\centering
\small
\caption{Top-3 features by single-feature $\auc$ ranking. Bold: shared across all continuation-mode models. Spearman $\rho$ (continuation pairs): Pythia--OLMo 0.95, Pythia--Qwen-cont 0.91, OLMo--Qwen-cont 0.85.}
\label{tab:top_features}
\resizebox{\columnwidth}{!}{%
\begin{tabular}{cllll}
\toprule
Rank & Pythia & OLMo & Qwen-rw & Qwen-cont \\
\midrule
1 & \textbf{\texttt{surp\_d1\_std}} & \textbf{\texttt{surp\_std}} & \texttt{freq\_entropy} & \textbf{\texttt{surp\_d1\_std}} \\
2 & \textbf{\texttt{surp\_d2\_std}} & \textbf{\texttt{surp\_d1\_std}} & \textbf{\texttt{surp\_std}} & \textbf{\texttt{surp\_std}} \\
3 & \textbf{\texttt{surp\_std}} & \textbf{\texttt{surp\_d2\_std}} & \textbf{\texttt{surp\_d1\_std}} & \textbf{\texttt{surp\_d2\_std}} \\
\bottomrule
\end{tabular}
}
\end{table}

%%% Table S1: Cross-Transfer Confusion Matrix (exp-013)
\begin{table*}[t]
  \centering
  \caption{Cross-model transfer accuracy matrix. Each cell reports the six-class depth classification accuracy of a detector trained on one model's outputs (row) and evaluated on another model's outputs (column). Diagonal entries (\textbf{bold}) indicate within-model performance. Transfer retention is computed as the ratio of off-diagonal accuracy to the geometric mean of the two corresponding diagonal entries.}
  \label{tab:s1_transfer}
  \begin{tabular}{lcccc}
    \toprule
    \multirow{2}{*}{\textbf{Train Model}} & \multicolumn{4}{c}{\textbf{Test Model}} \\
    \cmidrule(lr){2-5}
    & Qwen-rw & Pythia & OLMo & Qwen-cont \\
    \midrule
    Qwen-rw   & \textbf{0.4229} & 0.2590 & 0.2653 & 0.3466 \\
    Pythia    & 0.2718 & \textbf{0.3856} & 0.3531 & 0.3287 \\
    OLMo      & 0.2768 & 0.3450 & \textbf{0.3822} & 0.3244 \\
    Qwen-cont & 0.3440 & 0.3284 & 0.3112 & \textbf{0.3875} \\
    \bottomrule
  \end{tabular}
  \vspace{4pt}

  \footnotesize
  \textit{Transfer retention} between models $A$ and $B$ is defined as $\frac{1}{2}\bigl(\frac{\mathrm{acc}(A \to B)}{\sqrt{\mathrm{acc}(A \to A) \cdot \mathrm{acc}(B \to B)}} + \frac{\mathrm{acc}(B \to A)}{\sqrt{\mathrm{acc}(A \to A) \cdot \mathrm{acc}(B \to B)}}\bigr)$.
  Notable pairs: Pythia$\leftrightarrow$OLMo 91--92\%, Qwen-cont$\leftrightarrow$Pythia 83--85\%, Qwen-cont$\leftrightarrow$OLMo 80--85\%, Qwen-rw$\to$Pythia/OLMo 61--63\%.
\end{table*}

%%% Table S2: Decoding Ablation AUC (exp-003)
\begin{table*}[t]
  \centering
  \caption{Pairwise AUC with bootstrap 95\% CIs across depth pairs for four decoding configurations (all using Qwen2.5-1.5B-Instruct in rewrite mode; 10{,}000 bootstrap resamples). The horizontal rule between depth pairs 1\,vs.\,2 and 2\,vs.\,3 marks the empirical detectability boundary $\Kempirical=2$. Bottom rows report cross-configuration Spearman $\rho$ on per-feature discriminative-power rankings.}
  \label{tab:s2_decoding}
  \footnotesize
  \begin{tabular}{lcccc}
    \toprule
    \textbf{Depth Pair} & \textbf{Greedy} & \textbf{Nuc $p{=}0.9$, $T{=}1.0$} & \textbf{Nuc $p{=}0.95$, $T{=}0.7$} & \textbf{Nuc $p{=}0.95$, $T{=}1.2$} \\
    \midrule
    0 vs.\ 1 & $.980\;[.977,\,.982]$ & $.995\;[.993,\,.996]$ & $.988\;[.986,\,.990]$ & $.998\;[.998,\,.999]$ \\
    1 vs.\ 2 & $.689\;[.680,\,.698]$ & $.687\;[.677,\,.695]$ & $.689\;[.680,\,.697]$ & $.704\;[.695,\,.713]$ \\
    \cmidrule(lr){1-5}
    2 vs.\ 3 & $.561\;[.552,\,.571]$ & $.592\;[.582,\,.601]$\rlap{$^\dagger$} & $.590\;[.581,\,.600]$ & $.582\;[.572,\,.592]$ \\
    3 vs.\ 4 & $.540\;[.531,\,.550]$ & $.547\;[.537,\,.557]$ & $.539\;[.529,\,.549]$ & $.529\;[.519,\,.540]$ \\
    4 vs.\ 5 & $.514\;[.503,\,.524]$ & $.507\;[.496,\,.517]^{\ddagger}$ & $.535\;[.524,\,.545]$ & $.524\;[.513,\,.535]$ \\
    \midrule
    \multicolumn{5}{l}{\textbf{Cross-config feature Spearman $\rho$}} \\
    \midrule
    \multicolumn{2}{l}{Greedy vs.\ Nuc $p{=}0.9$} & 0.600 &
    \multicolumn{1}{l}{Nuc $p{=}0.9$ vs.\ $T{=}0.7$} & 0.782 \\
    \multicolumn{2}{l}{Greedy vs.\ $T{=}0.7$} & 0.914 &
    \multicolumn{1}{l}{Nuc $p{=}0.9$ vs.\ $T{=}1.2$} & 0.893 \\
    \multicolumn{2}{l}{Greedy vs.\ $T{=}1.2$} & 0.532 &
    \multicolumn{1}{l}{$T{=}0.7$ vs.\ $T{=}1.2$} & 0.668 \\
    \bottomrule
    \multicolumn{5}{l}{${}^{\dagger}$\,CI upper bound (0.601) marginally exceeds $\theta{=}0.60$ by 0.001.} \\
    \multicolumn{5}{l}{${}^{\ddagger}$\,Permutation $p = 0.12$ (not significant). All other pairs $p < 0.01$.}
  \end{tabular}
\end{table*}

%%% Table S2b: Extended K* CI (all models including new additions)
\begin{table*}[t]
  \centering
  \scriptsize
  \caption{Extended pairwise $\auc$ with bootstrap 95\% CIs for all model configurations including new additions.
  This table extends Table~\ref{tab:kstar_ci} in the main text with Qwen-base, LLaMA-8B, Mistral-7B, Pythia-6.9B (nucleus), and Gemma-2-2B configurations.
  $\Kempirical = 2$ holds for 6/10 conditions at $\theta{=}0.60$ (OLMo-1B, Mistral-7B, Pythia-6.9B, Gemma-2-2B: $\Kempirical{=}1$).}
  \label{tab:kstar_ci_extended}
  \resizebox{\textwidth}{!}{%
  \begin{tabular}{lcccccccccccccccccccc}
    \toprule
     & \multicolumn{2}{c}{Qwen-rw} & \multicolumn{2}{c}{Pythia} & \multicolumn{2}{c}{OLMo} & \multicolumn{2}{c}{Qwen-cont} & \multicolumn{2}{c}{Qwen-7B\rlap{$^\ddagger$}} & \multicolumn{2}{c}{Qwen-base} & \multicolumn{2}{c}{LLaMA-8B} & \multicolumn{2}{c}{Mistral-7B} & \multicolumn{2}{c}{Pythia-6.9B\rlap{$^\S$}} & \multicolumn{2}{c}{Gemma-2B} \\
    \cmidrule(lr){2-3}\cmidrule(lr){4-5}\cmidrule(lr){6-7}\cmidrule(lr){8-9}\cmidrule(lr){10-11}\cmidrule(lr){12-13}\cmidrule(lr){14-15}\cmidrule(lr){16-17}\cmidrule(lr){18-19}\cmidrule(lr){20-21}
     & $\auc$ & 95\% CI & $\auc$ & 95\% CI & $\auc$ & 95\% CI & $\auc$ & 95\% CI & $\auc$ & 95\% CI & $\auc$ & 95\% CI & $\auc$ & 95\% CI & $\auc$ & 95\% CI & $\auc$ & 95\% CI & $\auc$ & 95\% CI \\
    \midrule
    0v1 & .9956 & & .9932 & & .9860 & & .9971 & & .9638 & [.961,\,.967] & .987 & [.985,\,.989] & .9224 & [.917,\,.927] & .9842 & [.982,\,.986] & .9599 & & .953 & [.949,\,.957] \\
    1v2 & .6955 & [.686,\,.704]\rlap{$^*$} & .6539 & [.645,\,.663]\rlap{$^*$} & .5080 & [.496,\,.519] & .6112 & [.602,\,.621]\rlap{$^*$} & .6168 & [.608,\,.626]\rlap{$^*$} & .626 & [.617,\,.635]\rlap{$^*$} & .6123 & [.603,\,.622]\rlap{$^*$} & .5934 & [.584,\,.603] & .5955 & & .5981 & [.588,\,.608] \\
    2v3 & .5893 & [.580,\,.599] & .5514 & [.542,\,.563] & .4973 & [.486,\,.509] & .5609 & [.551,\,.571] & .5492 & [.539,\,.559] & .550 & [.540,\,.560] & .5273 & [.517,\,.538] & .5402 & [.530,\,.551] & .5298 & & .534 & [.523,\,.545] \\
    3v4 & .5435 & & .5152 & & .5024 & & .5409 & & .5276 & & .529 & [.519,\,.540] & .5168 & [.506,\,.528] & .5233 & [.513,\,.534] & .5135 & & .5033 & [.492,\,.514] \\
    4v5 & .5340 & & .5136 & & .4942 & & .5141 & & .5085 & & .518 & [.508,\,.529] & .5159 & [.505,\,.527] & .5180 & [.507,\,.529] & -- & -- & .5025 & [.492,\,.514] \\
    \bottomrule
    \multicolumn{21}{l}{\footnotesize $^*$Permutation $p < 0.001$ for all boundary comparisons.} \\
    \multicolumn{21}{l}{\footnotesize $^\ddagger$Generator: Qwen2.5-7B (fp16); scorer: Qwen2.5-1.5B. ``--'': experiment pending.} \\
    \multicolumn{21}{l}{\footnotesize $^\S$Pythia-6.9B evaluated separately; nucleus sampling ($p{=}0.95$, $T{=}1.0$), same as all other models.}
  \end{tabular}}% end resizebox
\end{table*}

%%% Table S3: LLM Zero-Shot Confusion Matrices (exp-011)
\begin{table*}[t]
  \centering
  \caption{LLM zero-shot depth classification: $5{\times}5$ confusion matrices
    (Exp-011). Both frontier LLMs achieve high binary detection accuracy
    (GPT-5.5: 0.8780; Claude Opus~4.6: 0.8818) but near-chance 5-class
    accuracy (${\sim}0.33$ vs.\ OBD-LightGBM baseline 0.386, evaluated on the full test set $N{=}5{,}000$). GPT-5.5
    strongly identifies depth-0 text (recall${=}0.7700$) but collapses on
    depth-4 (recall${=}0.0800$), misclassifying 50/100 depth-4 samples as
    depth-2. Claude Opus~4.6 shows more uniform errors across depths 1--4,
    with depth-4 recall (0.3700) substantially exceeding GPT-5.5.
    Per-depth recall comparison with OBD-LightGBM in Table~\ref{tab:s3c_perdepth}.}
  \label{tab:s3_llm_zeroshot}

  \vspace{0.5em}
  \begin{subtable}{0.48\textwidth}
    \centering
    \caption{GPT-5.5 ($N{=}500$, acc${=}0.332$, MAE${=}1.102$)}
    \label{tab:s3a_gpt55}
    \small
    \begin{tabular}{lcccccc}
      \toprule
      & \multicolumn{5}{c}{Predicted Depth} & \\
      \cmidrule(lr){2-6}
      True & 0 & 1 & 2 & 3 & 4 & Rec. \\
      \midrule
      0 & \textbf{77} & 18 &  4 &  0 &  1 & .770 \\
      1 &  9 & \textbf{19} & 36 & 24 & 12 & .190 \\
      2 &  8 & 10 & \textbf{38} & 35 &  9 & .380 \\
      3 & 10 & 16 & 35 & \textbf{24} & 15 & .240 \\
      4 & 11 & 13 & 50 & 18 &  \textbf{8} & .080 \\
      \bottomrule
    \end{tabular}
  \end{subtable}
  \hfill
  \begin{subtable}{0.48\textwidth}
    \centering
    \caption{Claude Opus~4.6 ($N{=}499$, acc${=}0.331$, MAE${=}1.150$)}
    \label{tab:s3b_claude}
    \small
    \begin{tabular}{lcccccc}
      \toprule
      & \multicolumn{5}{c}{Predicted Depth} & \\
      \cmidrule(lr){2-6}
      True & 0 & 1 & 2 & 3 & 4 & Rec. \\
      \midrule
      0 & \textbf{69} & 26 &  3 &  0 &  2 & .690 \\
      1 &  8 & \textbf{21} & 32 & 11 & 28 & .210 \\
      2 &  6 & 36 & \textbf{28} &  5 & 25 & .280 \\
      3 &  7 & 20 & 31 & \textbf{10} & 31 & .100 \\
      4 &  7 & 28 & 22 &  6 & \textbf{37} & .370 \\
      \bottomrule
    \end{tabular}
  \end{subtable}
\end{table*}

\begin{table}[t]
  \centering
  \small
  \caption{Per-depth recall (\%) on balanced $N{=}500$ test set. OBD's nominal overall accuracy gap originates from depth-0; all methods fail at depths $\geq 2$, consistent with $\Kempirical = 2$.}
  \label{tab:s3c_perdepth}
  \begin{tabular}{lccccc}
    \toprule
    & Depth 0 & Depth 1 & Depth 2 & Depth 3 & Depth 4 \\
    \midrule
    OBD-LightGBM & \textbf{96} & \textbf{39} & 14 & 19 & 22 \\
    GPT-5.5      & 77 & 19 & \textbf{38} & \textbf{24} &  8 \\
    Claude Opus~4.6 & 69 & 21 & 28 & 10 & \textbf{37} \\
    \bottomrule
  \end{tabular}
\end{table}

%%% Table: Token Length and EOS Truncation Analysis
\begin{table*}[t]
  \centering
  \small
  \caption{Mean subword token count per depth and EOS truncation rate (\%) for depth~1+ across all model configurations (generation limit: 256 tokens). Depth~0 (human text) is shared across configurations. Depths 1--5 show no systematic length trend within any model; the 0v1 length asymmetry (human texts shorter) makes detection \emph{more conservative}, not inflated. See Section~\ref{sec:length_confound} for discussion.}
  \label{tab:token_length}
  \begin{tabular}{lcccccccc}
    \toprule
    & \multicolumn{6}{c}{\textbf{Mean Token Count}} & & \\
    \cmidrule(lr){2-7}
    \textbf{Model} & Depth 0 & Depth 1 & Depth 2 & Depth 3 & Depth 4 & Depth 5 & \textbf{EOS Trunc.\ (\%)} & \textbf{Severity} \\
    \midrule
    Qwen-rw (1.5B)  & 191 & 254 & 252 & 254 & 254 & 254 & 7--11 & Mild \\
    Pythia-1.4B      & 192 & 214 & 215 & 216 & 217 & 218 & 34--36 & High \\
    OLMo-1B          & 192 & 256 & 256 & 256 & 256 & 256 & $<$1 & Negligible \\
    Qwen-cont (1.5B) & 191 & 244 & 238 & 240 & 240 & 238 & 17--20 & Moderate \\
    Qwen-7B          & 191 & 196 & 195 & 193 & 197 & 197 & 45--47 & High \\
    \midrule
    \multicolumn{9}{l}{\textit{Decoding ablation (Qwen-rw 1.5B)}} \\
    \quad Greedy          & 191 & 253 & 252 & 253 & 254 & 254 & 8--12 & Mild \\
    \quad Nuc $p{=}0.9$   & 191 & 254 & 251 & 254 & 253 & 254 & 8--16 & Mild \\
    \quad $T{=}0.7$       & 191 & 253 & 251 & 253 & 253 & 254 & 8--14 & Mild \\
    \quad $T{=}1.2$       & 191 & 254 & 252 & 253 & 254 & 254 & 8--16 & Mild \\
    \bottomrule
  \end{tabular}

  \vspace{4pt}
  \footnotesize
  EOS truncation rate: percentage of depth~1+ samples where generation terminated via EOS before reaching the 256-token limit. Ranges span depths 1--5.
\end{table*}

\clearpage
\section{Supplementary Details}

%%% Table S4: Dataset Statistics
\begin{table}[t]
  \centering
  \small
  \caption{Dataset statistics for the generational depth corpus. Depth~0 contains human-written seed texts sampled from the C4 English validation split and truncated to 256 subword tokens. Depths~1--5 are recursively generated by Qwen2.5-1.5B-Instruct (nucleus sampling, $p{=}0.95$, $T{=}1.0$). Token counts use whitespace tokenization. Statistics shown for the primary Qwen-rewrite corpus; the other three model-mode configurations (Pythia-cont, OLMo-cont, Qwen-cont) share the same generation pipeline and corpus structure.}
  \label{tab:s4_dataset_stats}
  \begin{tabular}{clrl}
    \toprule
    \textbf{Depth} & \textbf{$N$ samples} & \textbf{Mean tokens} & \textbf{Source} \\
    \midrule
    0 & 5{,}000 & 147.8 & C4 validation (human) \\
    1 & 5{,}000 & 207.1 & Qwen2.5-1.5B rewrite \\
    2 & 5{,}000 & 206.1 & Qwen2.5-1.5B rewrite \\
    3 & 5{,}000 & 208.0 & Qwen2.5-1.5B rewrite \\
    4 & 5{,}000 & 207.7 & Qwen2.5-1.5B rewrite \\
    5 & 5{,}000 & 207.7 & Qwen2.5-1.5B rewrite \\
    \midrule
    \multicolumn{2}{l}{Total} & \multicolumn{2}{l}{30{,}000 samples} \\
    \bottomrule
  \end{tabular}
\end{table}

%%% Table S5: Complete 15-D Feature Descriptions
\begin{table*}[t]
  \centering
  \caption{Complete description of the 15-dimensional feature set used by the OBD classifier. Features are grouped into three families: surprisal statistics (9~dims) computed from Qwen2.5-1.5B token-level log-probabilities, lexical diversity (3~dims), and frequency tail shape (3~dims). All features are computed per document.}
  \label{tab:s5_features}
  \begin{tabular}{llp{0.45\textwidth}l}
    \toprule
    \textbf{Group} & \textbf{Feature} & \textbf{Description} & \textbf{Range} \\
    \midrule
    \multirow{9}{*}{\rotatebox[origin=c]{90}{\textsc{Surprisal}}}
    & \texttt{surp\_mean}   & Mean token surprisal $\frac{1}{T}\sum_t -\log p(x_t \mid x_{<t})$ & $\mathbb{R}^+$ \\
    & \texttt{surp\_std}    & Standard deviation of token surprisal & $\mathbb{R}^+$ \\
    & \texttt{surp\_skew}   & Skewness of the surprisal distribution & $\mathbb{R}$ \\
    & \texttt{surp\_kurt}   & Excess kurtosis of the surprisal distribution & $\mathbb{R}$ \\
    & \texttt{surp\_d1\_mean} & Mean of first-order surprisal differences $\Delta s_t = s_{t+1} - s_t$ & $\mathbb{R}$ \\
    & \texttt{surp\_d1\_std}  & Std of first-order surprisal differences & $\mathbb{R}^+$ \\
    & \texttt{surp\_d1\_skew} & Skewness of first-order surprisal differences & $\mathbb{R}$ \\
    & \texttt{surp\_d2\_mean} & Mean of second-order differences $\Delta^2 s_t$ & $\mathbb{R}$ \\
    & \texttt{surp\_d2\_std}  & Std of second-order surprisal differences & $\mathbb{R}^+$ \\
    \midrule
    \multirow{3}{*}{\rotatebox[origin=c]{90}{\textsc{Lex}}}
    & \texttt{ttr}          & Type-token ratio: $|V| / T$ & $[0, 1]$ \\
    & \texttt{hapax\_ratio} & Hapax legomena ratio: freq-1 types / unique types & $[0, 1]$ \\
    & \texttt{self\_bleu}   & Mean 4-gram BLEU against 100 random same-depth samples & $[0, 1]$ \\
    \midrule
    \multirow{3}{*}{\rotatebox[origin=c]{90}{\textsc{Freq}}}
    & \texttt{freq\_kurtosis}  & Excess kurtosis of word frequency distribution & $\mathbb{R}$ \\
    & \texttt{freq\_entropy}   & Shannon entropy of word frequencies: $-\sum p_w \log_2 p_w$ & $\mathbb{R}^+$ \\
    & \texttt{low\_freq\_ratio} & Proportion of vocabulary with count $\leq 2$ & $[0, 1]$ \\
    \bottomrule
  \end{tabular}
\end{table*}

%%% Tables S6 & S7: Side by side in a single table* to avoid table*/table overlap
\begin{table*}[t]
\begin{minipage}[t]{0.45\textwidth}
  \centering
  \small
  \captionof{table}{LightGBM hyperparameters for the pairwise OBD (Ordinal Binary Decomposition) classifier used throughout the paper. One binary classifier is trained per adjacent depth pair $(k, k{+}1)$; 5-fold GroupKFold cross-validation is grouped by \texttt{doc\_id} to prevent data leakage from the same seed document appearing in both train and test.}
  \label{tab:s6_lgbm_params}
  \begin{tabular}{ll}
    \toprule
    \textbf{Parameter} & \textbf{Value} \\
    \midrule
    \texttt{n\_estimators}      & 200 \\
    \texttt{learning\_rate}     & 0.05 \\
    \texttt{max\_depth}         & 6 \\
    \texttt{num\_leaves}        & 31 \\
    \texttt{verbose}            & $-1$ \\
    \texttt{n\_jobs}            & $-1$ (all cores) \\
    Cross-validation            & 5-fold GroupKFold \\
    Grouping key                & \texttt{doc\_id} \\
    Objective                   & Binary (per depth pair) \\
    \bottomrule
  \end{tabular}
\end{minipage}\hfill
\begin{minipage}[t]{0.52\textwidth}
  \centering
  \small
  \captionof{table}{Neural probe architecture and training details (Exp-015). Applied to all three model families (Qwen2.5-1.5B, Pythia-1.4B, OLMo-1B): each model's frozen hidden states are average-pooled into document embeddings and classified by a two-hidden-layer MLP.}
  \label{tab:s7_neural_probe}
  \resizebox{\textwidth}{!}{%
  \begin{tabular}{ll}
    \toprule
    \textbf{Component} & \textbf{Details} \\
    \midrule
    \multicolumn{2}{l}{\textit{Embedding extraction}} \\
    Base models       & Qwen2.5-1.5B / Pythia-1.4B / OLMo-1B (frozen) \\
    Representation    & Last hidden state, attention-masked avg pool \\
    Embedding dim     & 1536 / 2048 / 2048 \\
    Max sequence len   & 512 tokens \\
    \midrule
    \multicolumn{2}{l}{\textit{MLP classifier}} \\
    Architecture      & $d_{\text{emb}} \to 512 \to 256 \to 6$ \\
    Activations       & ReLU \\
    Dropout           & 0.1 (after each hidden layer) \\
    Output            & 6-class softmax \\
    \midrule
    \multicolumn{2}{l}{\textit{Training}} \\
    Optimizer         & Adam ($\text{lr} = 10^{-3}$, weight decay $= 10^{-4}$) \\
    Loss              & Cross-entropy \\
    Batch size        & 256 \\
    Max epochs        & 20 \\
    Early stopping    & Patience = 5 (validation loss) \\
    Cross-validation  & 5-fold GroupKFold \\
    \bottomrule
  \end{tabular}%
  }
\end{minipage}
\end{table*}

%%% Table S8: Per-Model Bootstrap CI Results (exp-007)
\begin{table*}[t]
  \centering
  \footnotesize
  \caption{Pairwise AUC with bootstrap 95\% confidence intervals for four model-mode configurations (Exp-007, 1{,}000 bootstrap resamples). The horizontal rule separates depth pairs above and below the detectability threshold $\theta = 0.60$. Permutation $p$-values test $H_0\colon \auc = 0.50$. All pairs are significant at $p < 0.01$ except OLMo 1\,vs.\,2 through 4\,vs.\,5 (all $p > 0.09$), consistent with $\Kempirical{=}1$ for OLMo.}
  \label{tab:s8_cv_folds}
  \begin{tabular}{lcccc}
    \toprule
    \textbf{Depth Pair} & \textbf{Qwen-rewrite} & \textbf{Pythia-1.4B} & \textbf{OLMo-1B} & \textbf{Qwen-cont} \\
    \midrule
    0 vs.\ 1
      & $.996\;[.995,\,.997]$
      & $.993\;[.992,\,.994]$
      & $.986\;[.984,\,.988]$
      & $.997\;[.996,\,.998]$ \\
    1 vs.\ 2
      & $.696\;[.686,\,.704]$
      & $.654\;[.645,\,.663]$
      & $.508\;[.496,\,.519]$
      & $.611\;[.602,\,.621]$ \\
    \cmidrule(lr){1-5}
    2 vs.\ 3
      & $.589\;[.580,\,.599]$
      & $.551\;[.542,\,.563]$
      & $.497\;[.486,\,.509]$
      & $.561\;[.551,\,.571]$ \\
    3 vs.\ 4
      & $.544\;[.533,\,.555]$
      & $.515\;[.504,\,.526]$
      & $.502\;[.491,\,.513]$
      & $.541\;[.531,\,.551]$ \\
    4 vs.\ 5
      & $.534\;[.524,\,.545]$
      & $.514\;[.502,\,.525]$
      & $.494\;[.483,\,.506]$
      & $.514\;[.504,\,.526]$ \\
    \bottomrule
  \end{tabular}

  \vspace{4pt}
  \footnotesize
  Each cell: AUC $[95\%$ CI$]$. CI from 1{,}000 bootstrap resamples; $p$-values from 500-resample permutation test.
  All pairs $p < 0.01$ except OLMo 1v2--4v5 ($p > 0.09$).
\end{table*}

%%% Table S9: Binary Detector Spearman ρ (exp-012)
\begin{table*}[t]
  \centering
  \small
  \caption{Spearman rank correlation ($\rho$) between binary detector confidence scores and true generational depth, and 6-class depth classification accuracy (Exp-012). Three metrics: log-PPL from the reference model (DivEye), log-PPL from the observer model, and PPL ratio (Binoculars; observer: Pythia-1.4B, performer: OLMo-1B). Random baseline for 6-class accuracy: 0.167. While statistically significant, correlations are driven by the 0v1 boundary; confidence scores do not monotonically track depth beyond generation~1. Qwen-rewrite exhibits reversed sign ($\rho > 0$), indicating that generation mode alters the direction of binary confidence's relationship with depth.}
  \label{tab:s9_binary_spearman}
  \begin{tabular}{lcccc}
    \toprule
    & \textbf{Qwen-rw} & \textbf{Pythia} & \textbf{OLMo} & \textbf{Qwen-cont} \\
    \midrule
    \multicolumn{5}{l}{\textit{Spearman $\rho$ with true depth}} \\
    \quad log-PPL (ref)      & $-0.123$ & $-0.297$ & $-0.378$ & $-0.389$ \\
    \quad log-PPL (observer) & $+0.082$ & $-0.453$ & $-0.461$ & $-0.162$ \\
    \quad PPL ratio           & $+0.119$ & $-0.392$ & $-0.407$ & $-0.058$ \\
    \midrule
    \multicolumn{5}{l}{\textit{Classification}} \\
    \quad 6-class accuracy    & 0.199    & 0.130    & 0.127    & 0.183 \\
    \midrule
    \multicolumn{5}{l}{Random baseline: 0.167 \quad $\Kempirical^{\text{binary}} = 1$ \quad $\Kempirical^{\text{OBD}} = 2$} \\
    \bottomrule
  \end{tabular}
\end{table*}

%%% Table S10: Pythia Decoding Ablation (exp-020)
\begin{table*}[t]
  \centering
  \small
  \caption{Pairwise $\auc$ with bootstrap 95\% CIs across depth pairs for Pythia-1.4B under four decoding configurations (Exp-020; 10{,}000 bootstrap resamples). Complements Table~\ref{tab:s2_decoding} (Qwen-rewrite) with a second model family. $^\dagger$Greedy decoding triggers repetition collapse in Pythia-1.4B (4-gram repeat rate: $d_0{=}0.01 \to d_1{=}0.72$), inflating 0v1 and 1v2 AUC via lexical degeneracy rather than depth-discriminative features.}
  \label{tab:s10_pythia_decoding}
  \begin{tabular}{lcccc}
    \toprule
    \textbf{Depth Pair} & \textbf{Greedy}$^\dagger$ & \textbf{Nuc $p{=}0.9$, $T{=}1.0$} & \textbf{Nuc $p{=}0.95$, $T{=}0.7$} & \textbf{Nuc $p{=}0.95$, $T{=}1.2$} \\
    \midrule
    0 vs.\ 1 & $.999\;[.999,\,.999]$ & $.995\;[.994,\,.996]$ & $.996\;[.995,\,.997]$ & $.994\;[.993,\,.994]$ \\
    1 vs.\ 2 & $.884\;[.878,\,.890]$ & $.675\;[.667,\,.684]$ & $.728\;[.721,\,.736]$ & $.618\;[.608,\,.628]$ \\
    \cmidrule(lr){1-5}
    2 vs.\ 3 & $.531\;[.520,\,.541]$ & $.559\;[.549,\,.570]$ & $.598\;[.589,\,.606]$\rlap{$^\ddagger$} & $.531\;[.520,\,.541]$ \\
    3 vs.\ 4 & $.522\;[.511,\,.532]$ & $.524\;[.513,\,.534]$ & $.563\;[.553,\,.572]$ & $.520\;[.509,\,.531]$ \\
    4 vs.\ 5 & $.514\;[.504,\,.525]$ & $.501\;[.490,\,.512]$ & $.527\;[.518,\,.537]$ & $.510\;[.499,\,.521]$ \\
    \midrule
    $\Kempirical$ & 2 & 2 & 2 & 2 \\
    \bottomrule
    \multicolumn{5}{l}{${}^{\ddagger}$\,CI upper bound (0.606) exceeds $\theta{=}0.60$; point estimate (.598) below threshold.} \\
  \end{tabular}
\end{table*}

%%% Table S11: Length Ablation (exp-021)
\begin{table*}[t]
  \centering
  \small
  \caption{Pairwise $\auc$ across depth pairs for three segment lengths (Exp-021; Qwen2.5-1.5B base, continuation mode). $\Kempirical{=}2$ is confirmed for 128 and 256 tokens; 512 tokens yields point estimate $\Kempirical{=}2$ but the 1v2 bootstrap CI lower bound falls marginally below $\theta$.}
  \label{tab:s11_length_ablation}
  \begin{tabular}{lccc}
    \toprule
    \textbf{Depth Pair} & \textbf{128 tokens} & \textbf{256 tokens} & \textbf{512 tokens} \\
    \midrule
    0 vs.\ 1 & .9791 & .9914 & .9917 \\
    1 vs.\ 2 & .6219 & .6282 & .6073$^\dagger$ \\
    \cmidrule(lr){1-4}
    2 vs.\ 3 & .5291 & .5556 & .5721 \\
    3 vs.\ 4 & .5349 & .5298 & .5300 \\
    4 vs.\ 5 & .5012 & .5137 & .5176 \\
    \midrule
    $\Kempirical$ & 2 & 2 & 2$^\dagger$ \\
    \bottomrule
    \multicolumn{4}{l}{\footnotesize $^\dagger$CI $[0.5979,\,0.6166]$; lower bound ${<}\theta{=}0.60$.}
  \end{tabular}
\end{table*}

%%% Table S12: Domain Ablation
\begin{table*}[t]
  \centering
  \small
  \caption{Pairwise $\auc$ across depth pairs for three text domains (Qwen2.5-1.5B-Instruct, continuation mode). Tests generalization of $\Kempirical{=}2$ beyond C4 web text to structured and technical domains.}
  \label{tab:s12_domain_ablation}
  \begin{tabular}{lccc}
    \toprule
    \textbf{Depth Pair} & \textbf{C4 (web)} & \textbf{Wikipedia} & \textbf{arXiv} \\
    \midrule
    0 vs.\ 1 & .9971 & .9961 & .9948 \\
    1 vs.\ 2 & .6112 & .6193 & .8581 \\
    \cmidrule(lr){1-4}
    2 vs.\ 3 & .5609 & .5498 & .6436 \\
    3 vs.\ 4 & .5409 & .5231 & .5440 \\
    4 vs.\ 5 & .5141 & .5141 & .5299 \\
    \midrule
    $\Kempirical$ & 2 & 2 & 3 \\
    \bottomrule
  \end{tabular}
\end{table*}

%%% Table S13: Scorer Sensitivity Analysis
\begin{table*}[t]
  \centering
  \small
  \caption{Scorer sensitivity: $\Kempirical$ under self-scoring vs.\ cross-scoring. Self-scorer uses the generator's own model family for feature extraction; cross-scorer uses alternative models (Pythia-1.4B, OLMo-1B). $\Kempirical$ agreement across scorer conditions indicates the boundary is governed by the generation process. Cross-scoring reduces $\Kempirical$ only for the most marginal self-scorer condition (Qwen-continuation).}
  \label{tab:s13_scorer_sensitivity}
  \begin{tabular}{lccccc}
    \toprule
    & & \multicolumn{3}{c}{\textbf{Scorer}} & \\
    \cmidrule(lr){3-5}
    \textbf{Generator} & \textbf{Mode} & \textbf{Self} & \textbf{Pythia} & \textbf{OLMo} & \textbf{Agreement} \\
    \midrule
    Qwen-1.5B-Instruct & Rewrite      & 2 & 2 & 2 & \checkmark \\
    Qwen-1.5B-Instruct & Continuation & 2 & 1 & 1 & $\times$ \\
    Qwen-1.5B (base)   & Continuation & 2 & 2 & 2 & \checkmark \\
    Pythia-1.4B         & Continuation & 2 & 2 & 2 & \checkmark \\
    OLMo-1B             & Continuation & 1 & --- & --- & --- \\
    Qwen-7B             & Continuation & 2 & --- & --- & --- \\
    Mistral-7B          & Continuation & 1 & --- & --- & --- \\
    \bottomrule
  \end{tabular}
\end{table*}

\clearpage

\section{Extended Discussion}
\label{sec:extended_discussion}

\subsection{Theoretical Implications of $\Kempirical = 2$}

The DPI guarantees a finite $\Ktheory$; for general web text, $\Kempirical \leq 2$ across all 10 configurations (1B--8B, 6 families), while structured domains (arXiv) yield $\Kempirical{=}3$, aligning with model collapse convergence~\cite{shumailov2024collapse, dohmatob2024model, alemohammad2024selfconsuming}.
Decoding strategy modestly affects signal strength (1v2 $\auc \in [0.6866, 0.7039]$) but does not shift the boundary.

\subsection{Cross-Model Transfer Implications}

Cross-model classification retains 80--100\% accuracy with scorer-invariant feature rankings ($\rho > 0.90$), enabling single-classifier deployment across generators.
Absolute discriminability at the $\Kempirical$ boundary is scorer-sensitive: alternative scorers reduce Qwen-continuation $\Kempirical$ from 2 to 1, while rewrite-mode $\Kempirical = 2$ is preserved.
Model family similarity modulates transfer (Pythia$\leftrightarrow$OLMo: 90--92\% vs.\ Qwen-cont$\leftrightarrow$OLMo: 80--85\%).

\begin{minipage}{\columnwidth}
\centering
\captionof{table}{Cross-scorer validation on Qwen-generated text. Alternative scorers (Pythia-1.4B, OLMo-1B) preserve feature ranking ($\rho > 0.90$) and rewrite-mode $\Kempirical = 2$, but reduce continuation-mode $\Kempirical$ from 2 to 1.}
\label{tab:crossscorer}
\resizebox{\columnwidth}{!}{{\small
\begin{tabular}{llcccc}
\toprule
\textbf{Scorer} & \textbf{Gen.\ Data} & $\boldsymbol{K^*}$ & \textbf{AUC(0v1)} & \textbf{AUC(1v2)} & \textbf{Spearman}~$\boldsymbol{\rho}$ \\
\midrule
% --- Qwen-RW block ---
Qwen-1.5B (self) & Qwen-RW & 2 & --- & 0.6955 & --- \\
Pythia-1.4B       & Qwen-RW & 2 & 0.9923 & 0.7012 & 0.906 \\
OLMo-1B           & Qwen-RW & 2 & 0.9918 & 0.6947 & 0.920 \\
\midrule
% --- Qwen-cont block ---
Qwen-1.5B (self) & Qwen-cont & 2 & --- & 0.6112 & --- \\
Pythia-1.4B       & Qwen-cont & \textcolor{red}{1\textsuperscript{$\dagger$}} & 0.9900 & \underline{0.5923} & 0.942 \\
OLMo-1B           & Qwen-cont & \textcolor{red}{1\textsuperscript{$\dagger$}} & 0.9890 & \underline{0.5858} & 0.902 \\
\midrule
% --- Qwen-7B cross-scale block ---
Qwen-1.5B        & Qwen-7B-cont & 2 & 0.9644 & 0.6178 & --- \\
\bottomrule
\end{tabular}
}
}
\end{minipage}

\begin{table}[h]
\centering
\small
\caption{Three-way ANOVA (Generator $\times$ Scorer $\times$ Depth-Pair) on per-fold pairwise $\auc$ values. DV = per-fold AUC; $N{=}300$ (12 cells $\times$ 5 folds $\times$ 5 depth-pairs). Scorer is statistically significant but explains modest variance ($\eta^2_p{=}.166$) relative to Depth-Pair ($\eta^2_p{=}.997$), supporting near-scorer-independence of depth discrimination.}
\label{tab:anova_crossscorer}
\begin{tabular}{lrrrr}
\toprule
\textbf{Source} & \textbf{df} & $F$ & $p$ & $\eta^2_p$ \\
\midrule
Generator     & 2  & 851.91  & ${<}.001$  & .866 \\
Scorer        & 3  & 17.48  & ${<}.001$  & .166 \\
Depth-Pair    & 4  & 22737.70  & ${<}.001$  & .997 \\
Gen$\times$Scorer & 6 & 7.50 & ${<}.001$ & .146 \\
Gen$\times$Depth  & 8 & 913.41 & ${<}.001$ & .965 \\
Scorer$\times$Depth & 12 & 2.25 & $.010$ & .093 \\
Residual      & 264 & & & \\
\bottomrule
\end{tabular}
\end{table}

\subsection{Feature Importance Across Decoding Strategies}

Feature importance ordering varies across decoding strategies: greedy and $T{=}0.7$ share highly correlated rankings ($\rho = 0.914$), while greedy vs.\ $T{=}1.2$ yields lower correlation ($\rho = 0.53$, nonsignificant after Bonferroni).
The depth boundary is robust but the optimal feature subset is strategy-dependent.

\subsection{Scale and Neural Probe}
\label{sec:extended_neural_probe}

Cross-scale validation confirms $\Kempirical = 2$ from 1.5B through 8B: Qwen-7B (1v2 $\auc = 0.6168$ $[0.608, 0.626]$) and LLaMA-3.1-8B (1v2 $\auc = 0.6123$, 0v1 $= 0.9225$) both exceed $\theta$, with progressively lower 0v1 signal consistent with larger models generating more human-like text. Scaling to 13B+ remains open.

The neural probe yields model-dependent $\Kempirical^{\text{neural}}$: $3$ for Qwen (self-familiarity effect from instruction-tuned representations) but $2$ for Pythia; OLMo yields $\Kempirical^{\text{neural}} = 1$ (1v2 $\auc = 0.5054$), consistent with the 15D LightGBM boundary.
The 15D model-free features yield $\Kempirical = 2$ for Pythia and Qwen, while OLMo yields $\Kempirical = 1$ under the standardized pipeline.

\subsection{LLM Zero-Shot Baseline}
\label{sec:extended_llm}

GPT-5.5 achieves 33.2\% accuracy (MAE = 1.10) and Claude Opus~4.6 33.1\% (MAE = 1.15) on a balanced 500-sample test set, comparable to OBD-LightGBM's 38.0\% (95\% CI: $[0.340, 0.422]$; McNemar $p = 0.081$, $0.069$; Table~\ref{tab:s3_llm_zeroshot}).
Per-depth recall confirms convergent failure at depths $\geq 2$ (Section~\ref{sec:why_kstar}).
When 6-class probabilities are marginalized to 5 classes ($P(d{\geq}4)$), accuracy rises to 44.4\% but concentrates in depth-4 (recall: 22\%$\to$80\%) while depths 2--3 collapse to 2--8\%; we report raw 6-class results.

\subsection{Practical Considerations}

All filtering strategies degrade general capabilities (MMLU: 0.6018 $\to$ 0.5865--0.5941; HellaSwag: 0.6821 $\to$ 0.6638--0.6701); our 100\% synthetic setting is an upper bound; in mixed corpora, degradation scales with contamination fraction~\cite{shumailov2024collapse, dohmatob2024model}.
$\Kstar$-aware binary filtering reduces MMLU degradation by ${\sim}$50\% ($-$1.53\,pp $\to$ $-$0.77\,pp), requires no per-model calibration, and is combinable with distribution-matching methods (DSIR, DoReMi).

\subsection{Extended Limitations}
\label{sec:extended_limitations}

\emph{Prompt-chain versus retrain-chain.}
Retrain-chain validation (Pythia-1.4B LoRA, 5 generations) yields $\Kempirical{\geq}4$ (1v2 $\auc = 0.9484$, 2v3 $= 0.8888$, 3v4 $= 0.7982$), far exceeding prompt-chain $\Kempirical{=}2$.
The elevated signal arises from cumulative model-adaptation artifacts rather than text-intrinsic degradation; prompt-chain $\Kempirical$ is the operationally relevant boundary for data curation.

\emph{Threshold sensitivity.}
At $\theta = 0.55$, $\Kempirical = 3$ for Qwen-rewrite/continuation; at $\theta = 0.70$, $\Kempirical = 1$ across all configurations (Figure~\ref{fig:kstar_theta}).
Qwen-continuation's 1v2 CI lower bound exceeds $\theta$ by only 0.2\,pp; all ten boundary comparisons yield permutation $p < 0.001$.

\emph{Structural confound: token length.}
In continuation mode the generation prompt is discarded, so all depths $\geq 1$ consist entirely of machine-generated tokens and the 1v2 boundary compares texts differing only in conditioning context (human-authored vs.\ machine-generated prompt), not structural composition.
Depth~0 texts average ${\sim}$191 subword tokens (C4 truncation), while depth~1+ texts average 193--256 tokens depending on model (Table~\ref{tab:token_length}).
This asymmetry is directionally opposite to a length-inflation confound: shorter depth-0 texts provide less feature-extraction context, making 0v1 discrimination harder, not easier.
Within depths 1--5, mean token counts are stable (within-model range $\leq$ 4 tokens for all configurations except Qwen-continuation at 6 tokens), so the pairwise comparisons underlying $\Kempirical{=}2$ are unaffected by length.
EOS truncation rates are similarly depth-stable within each model: Qwen-rewrite 7--11\%, Pythia 34--36\%, OLMo $<$1\%, Qwen-continuation 17--20\%, Qwen-7B 45--47\%.
Models with the highest truncation (Pythia, Qwen-7B) still yield $\Kempirical{=}2$ with the same step-function AUC decay pattern as low-truncation models (Table~\ref{tab:kstar_ci}), confirming that EOS truncation does not drive the boundary.

\emph{Graduated filtering.}
All $\alpha$ values (0.1--0.9) underperform binary-ours on MMLU ($\alpha{=}0.1$: 0.5873; $\alpha{=}0.9$: 0.5892 vs.\ 0.5941); at $\alpha{=}0.9$, HellaSwag matches exactly (0.6701).
Binary-ours significantly outperforms all alternatives on MMLU ($p_{\text{corr}} \leq 0.021$); Graduated$_{0.9}$ is indistinguishable on HellaSwag ($p = 1.00$) but worse on MMLU ($p_{\text{corr}} = 0.011$).
\label{tab:pairwise_significance}

\subsection{Dose-Response: Contamination Ratio Effect}
\label{sec:dose_response}

The main filtering experiments (Section~\ref{sec:filtering}) use 100\% synthetic data to maximize effect size.
To characterize behavior at lower contamination ratios, we mix synthetic depths 1--5 with clean C4 at 10\%, 30\%, and 50\% contamination (Qwen-1.5B continuation, 5 seeds each; Table~\ref{tab:dose_response}).

\begin{table}[h]
\centering
\small
\caption{MMLU accuracy by contamination ratio (fraction of synthetic data in training set) and filtering strategy. Base model (no fine-tuning): 0.6018. Mean $\pm$ std over 5 seeds.}
\label{tab:dose_response}
\begin{tabular}{lccc}
\toprule
\textbf{Strategy} & \textbf{10\%} & \textbf{30\%} & \textbf{50\%} \\
\midrule
No-filter   & $0.6009 \pm .003$ & $0.5978 \pm .004$ & $0.5932 \pm .004$ \\
Binary-ours & $0.6011 \pm .003$ & $0.5988 \pm .003$ & $0.5956 \pm .004$ \\
Graduated   & $0.6008 \pm .003$ & $0.5982 \pm .004$ & $0.5944 \pm .004$ \\
\bottomrule
\end{tabular}
\end{table}

At 10\% contamination, no-filter MMLU ($0.6009$) remains within noise of the base model ($0.6018$, $\Delta = -0.09$\,pp), and all three strategies are indistinguishable ($\Delta_{\max} = 0.03$\,pp between strategies, $p > 0.8$).
At 30\%, differences emerge directionally ($\Delta = 0.10$\,pp, binary vs.\ no-filter) but remain non-significant ($p > 0.3$).
At 50\%, binary-ours shows a 0.24\,pp advantage over no-filter, approaching but not reaching significance ($p = 0.12$).
The strategy ordering (binary $>$ graduated $>$ no-filter) is consistent with the 100\% setting, but effect sizes scale linearly with contamination fraction.
This demonstrates that $\Kstar$-aware curation provides practical benefit primarily when contamination is severe enough to cause measurable quality degradation.

\subsection{Seed Robustness}
\label{sec:seed_robustness}

Because the random seed governs both corpus sampling and generation quality (not merely CV splits), absolute AUC values are seed-dependent.
We re-run three configurations with seed$\,{=}\,123$ (primary: seed$\,{=}\,42$; Table~\ref{tab:seed_robustness}).
$\Kempirical$ is stable in 2/3 cases: Pythia-1.4B on Wikipedia yields $\Kempirical{=}1$ under both seeds (1v2 $\Delta = -0.8$\,pp), and Qwen-1.5B on C4 yields $\Kempirical{=}2$ (1v2 $\Delta = +9.0$\,pp, both seeds above $\theta$).
Qwen-1.5B on Wikipedia shifts from $\Kempirical{=}2$ to $\Kempirical{=}3$ under seed$\,{=}\,123$, but the 2v3 AUC ($0.6067$) is marginal: one of five folds falls below $\theta$ (fold~5 $= 0.595$), so this does not constitute strong evidence for $\Kempirical{=}3$.
The ${\pm}9$\,pp absolute AUC variation across seeds strengthens the methodological case for threshold-based $\Kempirical$ over raw AUC comparisons: the qualitative boundary is robust even when point estimates shift.

\begin{table}[h]
\centering
\footnotesize
\caption{Seed robustness: $\Kempirical$ under two random seeds. 1v2 $\Delta$: change in 1v2 AUC from seed 42 to seed 123. $\Kempirical$ is stable in 2/3 configurations; the Qwen-Wiki shift to $\Kempirical{=}3$ is marginal (2v3 fold variance crosses $\theta$).}
\label{tab:seed_robustness}
\resizebox{\columnwidth}{!}{%
\begin{tabular}{llcccc}
\toprule
\textbf{Model} & \textbf{Domain} & \multicolumn{2}{c}{$\Kempirical$} & \textbf{1v2 $\Delta$} & \textbf{Stable?} \\
\cmidrule(lr){3-4}
 & & seed=42 & seed=123 & & \\
\midrule
Pythia-1.4B & Wikipedia & 1 & 1 & $-0.8$\,pp & \checkmark \\
Qwen-1.5B  & C4        & 2 & 2 & $+9.0$\,pp & \checkmark \\
Qwen-1.5B  & Wikipedia & 2 & 3$^\dagger$ & $+8.7$\,pp & $\times$ \\
\bottomrule
\multicolumn{6}{l}{\footnotesize $^\dagger$2v3 $\auc = 0.6067$; fold 5 $= 0.595 < \theta$. Marginal boundary.}
\end{tabular}}
\end{table}


\end{document}
